% =========================================================
% Proof: e-Characterization of Supremum (two proofs)
% Source: volume-iii/analysis/bounding/notes/notes-bounds-extremal-values.tex
% Mode: standard
% =========================================================

\subsection*{$\varepsilon$-Characterization of Supremum}
\label{prf:supremum-approximation}

\begin{remark*}[Return]
$\leftarrow$ \hyperref[prop:eps-char]{Back to Proposition in Notes}
\end{remark*}

\noindent\textit{Proof A — by contradiction.}

\begin{proof}
\Claim Let $A \subseteq \mathbb{R}$ be nonempty and bounded above, and let
$s = \sup A$.
Then for every $\varepsilon > 0$ there exists $a \in A$ with $s - \varepsilon < a$.

\Given $s = \sup A$; let $\varepsilon > 0$.
\Goal $\exists\, a \in A$ such that $s - \varepsilon < a$.

\Strategy Suppose for contradiction that no such $a$ exists.
Then $a \le s - \varepsilon$ for all $a \in A$, so $s - \varepsilon$ is an upper
bound of $A$.
But $s - \varepsilon < s$, contradicting $s$ being the \emph{least} upper bound.
\Therefore the assumption is false; such an $a$ exists. \AsReq
\end{proof}

\medskip
\noindent\textit{Proof B — direct.}

\begin{proof}
\Claim Same statement.

\Given $s = \sup A$; let $\varepsilon > 0$.
\Goal $\exists\, a \in A$ such that $s - \varepsilon < a$.

Since $s - \varepsilon < s$ and $s$ is the least upper bound, $s - \varepsilon$
is \emph{not} an upper bound of $A$.
\ByDef upper bound, there must exist $a \in A$ with $a > s - \varepsilon$. \AsReq
\end{proof}

\begin{remark*}[Proof shape]
Both proofs turn on the same fact: a value strictly below $\sup A$ cannot
be an upper bound.
Proof A reaches this via contradiction; Proof B reads it directly from the
definition of least upper bound.
The direct proof is shorter; the contradiction proof makes the logical
negation explicit, which can aid comprehension.
\end{remark*}

\begin{remark*}[Dependencies]
Definition of supremum (least upper bound).
\end{remark*}
